\section{The remaining small-field budget inequalities}
\label{sec:small-field-budget-selection}
The nonlinear step uses several estimates at different small powers
of the coupling. We verify their numerical compatibility before
choosing the stage thresholds. Throughout this section the block
factor and stage data are fixed. Write $z=\beta_F(x)$ for the shifted
fine coupling, $B=B(x)$ for the block radius, and $\tau=\tau(x)$ for
the fluctuation threshold. Thus $z\to\infty$, $B$ has upper exponent
zero, and $\tau$ has upper exponent $-1/2$. We also use the
stronger growth facts $B/\log z\to\infty$ and
$P_F(x)^2/\log z\to\infty$ for the actual radius and fine monitor.
They follow from their prescribed logarithmic growth laws; upper
exponent zero alone would not suffice for the exponential estimates.
All constants may depend
on the fixed derivative order, spatial weights, path prices and
positive support-width reserves.

\subsection{The low-degree action and its support tail}
The explicit low-degree correction has four contributions: the
quadratic reset term, cubic and quartic bare-action terms, and a
quadratic loop term. Their rooted jet prices, with the fixed support
weights already included, have the form
\begin{equation}\label{eq:small-field-four-costs}
       u\le C_2|\alpha|\tau^2
                  +C_3|\beta|\tau^3
                  +C_4|\beta|\tau^4+C_{\rm loop}\Lambda\tau^2.
\end{equation}
The constants include the finite coordinate and orientation costs,
and the anchored support prices for the corresponding degrees.
This is the direct sum of the quadratic and homogeneous action
bounds: a degree-$r$ term tested in directions of size $\tau$
contributes $\tau^r$, while its coefficient contributes the displayed
$|\alpha|$, $|\beta|$, or $\Lambda$. The stage fixes the derivative
order, spatial power and support exponent, so these constants are
unchanged when its actual threshold varies.

Assume $|\alpha|\le A$, $|\beta|\le2z$, and
$0\le\Lambda\le\Lambda_0(x)$, where $A$ is fixed and $\Lambda_0$
has upper exponent zero. The four terms in
\eqref{eq:small-field-four-costs} have upper exponents respectively
$-1,-1/2,-1,-1$. Hence one common nonnegative majorant for $u$
has upper exponent $-1/2$. Since $-1/2<-7/16$, eventually
\begin{equation}\label{eq:small-field-action-price}
                              u\le z^{-7/16}.
\end{equation}
If $\sigma>0$ is the fixed support width, then
$e^{-\sigma B/4}$ has every negative upper exponent. Multiplying the
same majorant by this tail gives, after increasing the threshold,
\begin{equation}\label{eq:small-field-action-tail}
                  e^{-\sigma B/4}u\le z^{-23/20}.
\end{equation}
Also $(a+3)/B\le\sigma/4$ eventually for any fixed target support
exponent $a$. The thresholds here are uniform over the stated ranges
of $\alpha,\beta,\Lambda$ and over the actual finite torus.

The reset bound itself follows from the existing coupling estimate.
If the running reset satisfies
$|\alpha+c_G\log L|\le A_0$, then
$|\alpha|\le A_0+|c_G\log L|$. Enlarging the right side by one
makes a fixed positive bound available for this chosen block factor.
It does not require selecting a different reset at each application.

\subsection{The strengthened cluster remainder}
Let $D=z^{3/8}$ be the extra-degree weight, so $D\ge1$ eventually. Write
$\mathcal P_\eta(x)=C_j(x)e^\eta\eta^{-pj}$ for the derivative price at
positive reserve $\eta$, where $j,p\in\mathbb N$ are the fixed derivative
order and spatial power. The factor $C_j(x)$ includes $(pj)!$ and the
lattice scale factor; it is common to all reserves but may vary with
the block radius. It has upper exponent zero. We suppress $x$ in these
prices below.
Suppose the ordinary component/cluster budget at reserve $i>0$
has already given
\[
            \mathcal P_i\,S(1-1/2)^{-1}\le z^q,
                         \qquad q\le-7/16.
\]
Here $S=(1+2^jB_1)\omega\ge0$ is the unstrengthened root price for
the second expansion, in the notation of
Section~\ref{sec:high-degree-cluster-prices}; it includes the component
derivative factor. The reserve $i$ lies
below the fixed available reserve $y>0$, and let $a\ge0$ be the target
support exponent. Retain the graded elementary
bound and the two smallness conditions from
Section~\ref{sec:high-degree-cluster-prices}, namely
\[
 \frac{4(2B+1)^4D r_0}{\delta_0}\le\frac12,\qquad
 \frac{4(2B+1)^4D S}{y}\le\frac12,
\]
with the available support exponent covering the reserve $a+2y$.
Here $r_0\ge0$ and $\delta_0>0$ are the first expansion's root price
and reserve. Also retain its seed-mark and footprint purchases, for
$M\ge0$, $\epsilon>0$ and $T\ge0$,
$D^3\le e^M$ and $D^3\le e^{\epsilon T}$, in the notation of that
section. These conditions concern the same incoming representation
and graded component prices. The extra-degree argument
of Section~\ref{sec:high-degree-cluster-prices} gives a residual
bound
\[
               D^{-3}\mathcal P_y\,D S(1-1/2)^{-1}.
\]
Since $i\le y$, the inverse powers satisfy
$y^{-pj}\le i^{-pj}$, while $e^i\ge1$. Consequently
$\mathcal P_y\le e^y\mathcal P_i$. The displayed remainder is at
most
\[
                    e^yD^{-2}z^q=e^y z^{q-3/4}.
\]
Now $q-3/4\le-19/16<-23/20$. The strict gap absorbs the fixed
factor $e^y$ and proves the common residual budget $z^{-23/20}$.
This calculation consumes the unstrengthened price $S$ built from
the original component price and its derivative factor; it does not
assume an independent price for the strengthened expansion.

\subsection{Polynomial localization costs and the incoming allowance}
The auxiliary periods in the five-period construction have side
lengths bounded by a fixed polynomial in $B$. Their volumes, local
ball counts and ordered-derivative transport costs therefore all
have upper exponent zero. Multiplication by any of them preserves
the upper exponent of a price.

In particular, localizing the ordinary error of size $z^{-7/16}$
produces a total auxiliary error of upper exponent $-7/16$. A further
fixed support conversion, including the constant exponential of the
support-cover bound, still has that exponent. Since
$-7/16<-1/8$, its final price is eventually at most $z^{-1/8}$.
Similarly, localizing a price $z^{-23/20}$, including the factor
$j!3^j$ from Section~\ref{sec:healthy-boundary-jets}, gives a price
of upper exponent $-23/20$.

Let $G(x)$ be the sum of the latter localized price and the finitely
many unlocalized $z^{-23/20}$ remainders, each with its fixed
coefficient. Then $G$ also has upper exponent $-23/20$. If $c>0$
is the fixed regular allowance constant, the strict inequality
$-23/20<-11/10$ gives
\[
                   G(x)\le\tfrac c2 z^{-11/10}
                             \le\tfrac c2 x^{-11/10}
\]
eventually; the second inequality uses $z\ge x>0$. Thus these
remainders use at most half of the available regular allowance.
All polynomial localization factors are paid before this final
comparison.

\subsection{Absorbing the forced correction at its prescribed price}
A second exponential absorption must retain its prescribed positive
strong forced price $v$. Let $a,b\ge0$ be its hard and soft input prices, with
\[
                       a\le v/8,\qquad b\le z^{-1/8}.
\]
Put $p=z^{1/16}$. The constructed strong price has the form
$v=C\exp[-r(\tau/g(\beta))^2]$ with fixed $C,r>0$, where
$g(\beta)=\beta^{-1/2}$. It therefore has the monitor estimate
$v\le C\exp(-cP_F(x)^2)$ uniformly for the allowed coupling band,
with fixed $C,c>0$. Thus $vp\le1$ eventually. Also, uniformly for
$b$ in the stated range,
\[
 b\le1,\qquad pb\le1/8,\qquad
        (2B+1)^4\{(8p)^{-1}+2b\}\le1
\]
eventually: the two powers in the last braces are strictly negative
and the prefactor has upper exponent zero.

Choose
\[
                \kappa=-\log(vp)\ge0,
                   \qquad I=e^\kappa a+2b.
\]
Then $I\le(8p)^{-1}+2b$, so $(2B+1)^4I\le1$ and in particular
$0\le I\le1$. Using $e^I-1\le2I$ gives
\[
 e^{-\kappa}(e^I-1)
   \le 2vp\{(8p)^{-1}+2b\}
   =v(1/4+4pb)\le3v/4\le v.
\]
These are the exact smallness and output inequalities required by
the forced-gas absorption, at the original allowance $v$. Its
uniform monitor estimate follows from the coupling-band comparison
for the gap. Indeed $\tau=z^{-1/2}P_F(x)$ and $\beta\ge z/4$ give
$\beta\tau^2\ge P_F(x)^2/4$. The fixed positive gap and normalization
constants are absorbed into $c$. All constants and the threshold are
chosen before the particular coupling in the band.

Each assertion above is eventually valid for a fixed transition.
Take the finite maximum of its required thresholds at each stage,
then enlarge the entry coupling accordingly. Auxiliary periods are
also chosen above their fixed geometric lower bounds, and the
volume growth exponent is increased to at least one. These choices
are preserved by any later increase of the delays. They supply the
numerical budgets; the actual local action identities, stationary
comparison and funded incoming representations are still required
where each budget is used.
